\section{构建空间对齐的多模态癌症数据}

模型能力的上限先受数据工程限制。Noetik 的路线不是只下载公共图片，而是在同一患者组织上执行多种测量，并尽量保留空间对应关系。本章沿着样本制备、成像、配准和数据层级，解释为什么这种数据既昂贵又适合多模态建模。

\subsection{从组织样本到数字切片}

数据流水线从真实肿瘤组织开始。每一步都可能引入批次效应、形变、信号漂移或缺失，因此后续模型不能把所有差异都解释为生物机制。为了理解最终 token 的来源，下面沿着组织处理、微阵列、H\&E、免疫荧光和数据平台依次查看每个物理步骤怎样留下统计痕迹。

\lecturefigure{slide-097-tissue-processing.jpg}{组织样本进入自动化处理与切片流程}{00:25:22}

图中嵌入视频展示实验室处理，而右侧保留讲座时点的数据计数。样本必须固定、切片、染色并数字化；同一组织在不同 assay 中可能经历不同物理处理。若配准不准，模型看到的“同一位置”可能实际上来自不同细胞群。

\lecturefigure{slide-104-tissue-microarray.jpg}{组织微阵列：在载玻片上组织大量小型 tissue core}{00:25:31}

Tissue microarray 把来自不同患者或区域的小圆柱状组织核心排列在同一载玻片上，便于批量染色和扫描。它提高吞吐量，但每个 core 只覆盖肿瘤的一小部分，不能代表整个患者。多 core 设计是在成本与空间异质性之间折中。

\lecturefigure{slide-108-he-tissue-cores.jpg}{H\&E 组织核心：低成本、广覆盖的形态学入口}{00:25:36}

H\&E 是 hematoxylin and eosin（苏木精--伊红）染色。苏木精主要让细胞核呈蓝紫色，伊红让胞质和基质呈粉红色。它廉价、临床常规且历史数据丰富，但只能间接反映细胞类型和分子状态，不能读取所有蛋白或基因表达。

\lecturefigure{slide-119-multimodal-core-grid.jpg}{同一组织核心在多种染色与成像通道下形成配对数据}{00:25:49}

这一网格显示多模态数据真正的价值：不是随机找一张 H\&E 和一张免疫荧光，而是尽可能让不同测量对应同一患者、相近区域。这样的配对支持 translation；模态独有信号又支持 disambiguation。配准误差、切片间差异和组织损耗必须进入误差预算。

\lecturefigure{slide-125-data-platform-ui.jpg}{数据平台将多患者、多核心和多通道组织成可浏览对象}{00:25:57}

界面截图说明数据不是一张“超大图”，而是层级数据库：患者包含切片，切片包含 core，core 包含空间位置，位置又有图像、蛋白和基因读数。训练切分必须在患者层完成；若同一患者的相邻 core 分到训练和测试，指标会因患者特异纹理泄漏而虚高。

\lecturefigure{slide-129-immunofluorescence.jpg}{多重免疫荧光显示细胞类型与蛋白标志物的空间组合}{00:26:04}

Multiplex immunofluorescence（多重免疫荧光）用多种带荧光标签的抗体检测蛋白标志物。颜色可帮助识别 T 细胞、B 细胞、巨噬细胞或肿瘤区域，但颜色只是可视化；模型输入通常是各通道的连续强度和分割结果。抗体特异性与成像饱和会限制解释。

\begin{warningbox}{第一类数据泄漏：患者与切片不能混切}
空间相邻 patch 极其相似。若按 patch 随机切分，模型可能记住患者、扫描仪或染色批次，而不是学习可泛化生物规律。评估应优先采用 held-out patient，并单独报告跨批次与跨机构表现。
\end{warningbox}

\subsection{四种模态与空间转录组}

课堂随后把数据归纳为四类：H\&E、蛋白成像、空间转录组和遗传测序。它们观察的是同一系统的不同层面，成本、分辨率与生物含义均不相同。这里的核心问题不是把四种文件格式统一，而是维护 patient--core--cell--gene 的层级索引，并知道每种测量在哪个空间与生物尺度上有效。

\lecturefigure{slide-134-multimodal-dataset.jpg}{四类模态：H\&E、蛋白、空间转录组与遗传测序}{00:29:19}

H\&E 提供形态，蛋白面板提供细胞表型，spatial transcriptomics（空间转录组）测量带坐标的 RNA 表达，genetic sequencing（遗传测序）提供突变等信息。前两者像图像，后两者更像稀疏高维计数和患者级属性；统一 token 化必须尊重这种差异。

\begin{knowledgebox}{第一次出现：Spatial transcriptomics}
空间转录组同时测量“某些基因表达了多少”和“这些表达发生在组织的哪里”。与普通 bulk RNA 相比，它保留空间；与单细胞 RNA 相比，它的测量单元可能覆盖多个细胞，分辨率和基因覆盖取决于 assay。
\end{knowledgebox}

\lecturefigure{slide-135-spatial-transcriptomics.jpg}{空间转录组全景：组织形态上叠加大量细胞与表达位置}{00:29:42}

远看是一块组织，近看却包含数量巨大的点位和细胞。读图要携带三个层级：组织结构决定邻域，细胞类型影响基线表达，局部信号通路改变状态。若把所有细胞打散为无序表格，会丢失免疫细胞是否进入肿瘤、是否围绕坏死区聚集等关键信息。

\lecturefigure{slide-153-live-zoom.jpg}{实时缩放到细胞级：宏观纹理由海量局部测量组成}{00:30:47}

实时 demo 的意义不是展示缩放技术，而是让“序列长度”具象化。一个 core 可含成千上万细胞，每个细胞又有上千基因计数；若全部展开为 token，自注意力的二次复杂度无法直接承受。系统必须采样、分块、压缩或采用层级注意力。

\lecturefigure{slide-175-gene-overlay.jpg}{在组织空间上叠加基因表达与细胞身份}{00:31:16}

叠加后的标签帮助研究者把分子信号映射回空间生态。这里容易出现因果错觉：某基因在某区域高表达，可能因为该区域细胞组成不同，而不一定是同一种细胞发生状态变化。模型需要同时控制 cell type、neighborhood 和 patient context。

\lecturefigure{slide-177-patient-cores-gene-table.jpg}{每位患者有多个 core，每个细胞对应高维基因读数}{00:32:09}

这张摘要图把数据张量化：患者 $p$ 下有 core $c$，core 中有细胞 $i$，每个细胞含基因计数向量 $x_{p,c,i,g}$。层级结构意味着有效样本量随问题而变：预测某个基因的局部分布可利用许多细胞，预测新患者治疗响应仍主要受患者数限制。

\begin{importantbox}{数据层级决定统计单位}
不能用“有数十亿 token”掩盖“只有约千名患者”。细胞级任务与患者级任务的独立样本数不同；任何声称临床泛化的指标都应明确切分单位、患者数量和置信区间。
\end{importantbox}

\teachervoice{00:25:23--00:32:10，老师用整段实验室视频强调数据集不是现成资产：组织制备、扫描、配准和多 assay 都是系统的一部分。随后通过不断缩放让听众看到 patient--core--cell--gene 的规模爆炸。}

\subsection{本章小结}

多模态癌症建模首先是数据对齐问题。H\&E 廉价而间接，蛋白成像更接近细胞表型，空间转录组保留位置与表达，遗传测序提供患者级变异。它们的组合支持翻译与消歧，但也引入配准误差、批次效应、患者泄漏和层级样本量错觉。

